\chapter{Introducción}

La creación de un club ambiental en la Universidad Politécnica de Sinaloa sirve como una estrategia más para potenciar la formación de una comunidad estudiantil con preocupación por la sustentabilidad y el cuidado ambiental. 

La implementación de este club ofrece amplios beneficios. Por un lado, la participación en este club ofrece a la comunidad estudiantil un medio para el desarrollo de valores y habilidades de gran valor. Por otro lado, genera para la institución un canal por el cual hacer de la conciencia ambiental un valor importante en la institución. Por último, aporta  a las comunidades civiles y estudiantiles del municipio  un nuevo medio por el cual obtener conocimiento acerca de sostenibilidad.

\section{Justificación}

El cuidado ambiental es un área importante en los planes de estudio de diversas licenciaturas de la Universidad Politécnica de Sinaloa, por lo que la promoción de actividades de cuidado ambiental es también una oportunidad de aprendizaje extra para los alumnos, además de que les proporciona una herramienta más para aplicar los conocimientos con los que ya cuentan. Además de esto, las actividades de voluntariado sirven a los alumnos como experiencias sociales positivas que forjan un sentido de comunidad y colectividad entre los alumnos, así como una oportunidad más de desarrollo de habilidades sociales.

\section{Objetivos}

\subsection{General}

Crear un programa general de voluntariado para la realización de actividades relacionadas al cuidado del medio ambiente tanto en la comunidad estudiantil de la Universidad Politécnica de Sinaloa como en comunidades fuera de la universidad, ya sea estudiantiles o civiles.

\subsection{Específicos}

\begin{itemize}
    \item Establecer un programa de compostaje comunitario en la Universidad Politécnica de Sinaloa. 
    \item Formar un equipo de alumnos voluntarios que realicen actividades de difusión y educación ambiental, tanto dentro como fuera de la universidad. 
    \item Colaborar con organizaciones ciudadanas para la promoción entre la comunidad estudiantil de actividades de cuidado ambiental en el municipio. 
\end{itemize}

A continuacin se muestra una lista de programas planeados para el club:

    \begin{table}[H]

\centering
\caption{Lista de programas a cargo del club ambiental}
\label{tab:porgramas_del_club_ambiental}

\begin{threeparttable}

{\small\setstretch{1.25}\rowcolors{2}{Grey}{White}

\begin{tabular}{M{0.15\linewidth}M{0.375\linewidth}M{0.375\linewidth}}\hline
    
    \theader{Programa} & \theader{Funcionamiento} & \theader{Objetivos} \\ \hline

    Compostaje & Compostaje de los residuos orgánicos universitarios y de comercios cercanos, realizado por alumnos voluntarios y utilizado como abono en áreas verdes universitarias y prácticas académicas.
    &
    \begin{itemize}
        \item Crear una herramienta más para el aprendizaje de los alumnos
        \item Fomentar el manejo de residuos responsable
    \end{itemize} \\

    Difusión & Actividades dentro y fuera de la universidad para dar difusión a temas de cuidado ambiental.
    &
    \begin{itemize}
        \item Crear conciencia de la importancia del cuidado ambiental
        \item Dar herramientas a los estudiantes para llevar una vida más ecológica
        \item Dar a conocer avances en tecnologías de cuidado ambiental
    \end{itemize} \\

    Limpiezas & Eventos que convoquen a la comunidad estudiantil para realizar limpieza, tanto dentro como fuera de la universidad.
    &
    \begin{itemize}
        \item Fomentar la disposición responsable de residuos
        \item Mantener una universidad libre de contaminación
        \item Recolectar materia compostable
    \end{itemize} \\

    Vinculación & Contacto con organizaciones locales para asistir a sus eventos como representantes universitarios y dar difusión de estos entre la comunidad estudiantil.
    &
    \begin{itemize}
        \item Dar apoyo a organizaciones que fomenten el cuidado ambiental
        \item Reforestación (a futuro)
        \item Cultivo de árboles fertilizados con el programa de compostaje y su uso para proyectos de reforestación universitarios o municipales
        \item Dar apoyo a la formación de un municipio más sustentable
    \end{itemize} \\

    \hline
\end{tabular}

}

\begin{tablenotes}\tiny % ex: \tnote{1} ... \item[1]
    \item[1] 
\end{tablenotes}

\end{threeparttable}
\end{table}

